\documentclass[11pt]{article}
\usepackage[a4paper,margin=1in]{geometry}
\usepackage{fontspec}
\usepackage[boldfont=false]{xeCJK}
\setCJKmainfont{FandolSong-Regular.otf}
\usepackage{amsmath}
\usepackage{amssymb}
\usepackage{array}
\usepackage{booktabs}
\usepackage{graphicx}
\usepackage{adjustbox}
\usepackage{enumitem}
\setlist[itemize]{topsep=2pt,partopsep=0pt,parsep=0pt,itemsep=2pt,leftmargin=1.4em}
\setlist[enumerate]{topsep=2pt,partopsep=0pt,parsep=0pt,itemsep=2pt,leftmargin=1.6em}
\usepackage{parskip}
\usepackage{fvextra}
\fvset{breaklines=true,breakanywhere=true,fontsize=\small}
\setlength{\parindent}{0pt}

\begin{document}

\begin{center}
  {\LARGE\bfseries Diseases and immunity}\\[0.35em]
  {\large Igcse Biology}
\end{center}
\vspace{0.6em}

\section*{Pathogens and disease}

A \textbf{pathogen} 病原体 is an organism that causes \textbf{disease} 疾病. A \textbf{transmissible disease} 传染病 is one in which the pathogen can pass from one \textbf{host} 宿主 to another.

\section*{How pathogens spread}

A pathogen is \textbf{transmitted} 传播 in two main ways:

\begin{itemize}
\item by \textbf{direct contact} 直接接触 — for example through blood and other \textbf{body fluids} 体液.

\item indirectly — for example from \textbf{contaminated} 污染 surfaces, food, animals, or the air.

\end{itemize}

\section*{The body's defences}

Your body has several \textbf{defences} 防御 that keep pathogens out or destroy them:

\begin{itemize}
\item \textbf{skin} — a barrier that covers and protects the body.

\item \textbf{hairs in the nose} — trap dust and pathogens in the air you breathe in.

\item \textbf{mucus} 黏液 — sticky liquid in the airways that traps pathogens.

\item \textbf{stomach acid} 胃酸 — kills most pathogens in your food.

\item \textbf{white blood cells} 白细胞 — find and destroy any pathogens that get inside.

\end{itemize}

\section*{Controlling the spread of disease}

Good public health stops disease spreading:

\begin{itemize}
\item a clean water supply, so drinking water carries no pathogens.

\item \textbf{hygienic} 卫生 food preparation and good personal hygiene, such as washing your hands.

\item proper waste disposal and \textbf{sewage treatment} 污水处理, so waste does not contaminate water or food.

\end{itemize}

\section*{Immunity (Supplement)}

\subsection*{Antigens and antibodies}

Every pathogen carries \textbf{antigens} 抗原 on its surface, and each kind of antigen has its own special shape.

\textbf{Antibodies} 抗体 are \textbf{proteins} 蛋白质 made by \textbf{lymphocytes} 淋巴细胞 (a kind of white blood cell). An antibody has a shape that is \textbf{complementary} 互补 to one antigen, so it fits only that antigen. When antibodies bind to a pathogen's antigens, they either destroy the pathogen directly or mark it so that \textbf{phagocytes} 吞噬细胞 destroy it.

\subsection*{Active immunity and vaccination}

\textbf{Active immunity} 主动免疫 is protection made by your \textbf{own} body producing antibodies. You gain it after an infection, or after \textbf{vaccination} 疫苗接种.

Vaccination works like this:

\begin{enumerate}
\item weakened pathogens, or just their antigens, are put into the body; they cannot make you ill.

\item the antigens make your lymphocytes produce antibodies.

\item the body also makes \textbf{memory cells} 记忆细胞 that stay for years.

\end{enumerate}

If the real pathogen enters later, the memory cells make antibodies very fast, so you do not become ill. This gives long-term protection. If most people in a group are vaccinated, the disease cannot spread easily.

\subsection*{Passive immunity}

\textbf{Passive immunity} 被动免疫 is protection from antibodies made by \textbf{another} body, not your own. For example, a baby receives antibodies from its mother across the \textbf{placenta} 胎盘 and in \textbf{breast milk} 母乳. This is why \textbf{breast-feeding} 母乳喂养 helps protect \textbf{infants} 婴儿. Passive immunity is only short-term, because \textbf{no} memory cells are made.

\section*{Cholera (Supplement)}

\textbf{Cholera} 霍乱 is a disease caused by a \textbf{bacterium} 细菌 that spreads in contaminated water. The cholera bacterium makes a \textbf{toxin} 毒素. This toxin causes \textbf{chloride ions} 氯离子 to be pumped into the \textbf{small intestine} 小肠. Water then moves into the intestine by \textbf{osmosis} 渗透. The result is \textbf{diarrhoea} 腹泻, \textbf{dehydration} 脱水 (loss of water) and a loss of \textbf{ions} 离子 from the blood.

\section*{Exam tips}

\begin{itemize}
\item Pathogen = organism that causes disease. Transmissible = can pass between hosts, by direct contact or indirectly.

\item Learn the five body defences: skin, nose hairs, mucus, stomach acid, white blood cells.

\item An antibody's shape is \textbf{complementary} to one antigen (like an enzyme and its substrate).

\item Active immunity (from infection or vaccination) makes \textbf{memory cells} → long-term. Passive immunity (placenta, breast milk) makes \textbf{no} memory cells → short-term.

\item Cholera: toxin → chloride ions into the gut → water in by osmosis → diarrhoea and dehydration.

\end{itemize}


\end{document}
